% Related lecture: SC4001-L06
% Sources: 6_cnn pp. 61--65 and 93--99; no formal-exam question reproduced
% Human verified: Yes for stated dimensions and example 1 numerics

\exercisegroup{SC4001-L06-EX}{第 06 讲：Convolutional Neural Networks}

\exercisemeta
  {SC4001-L06-EX}
  {2026-09-18}
  {6\_cnn pp.~61--65、93--99；例题数值与维度已独立核算}
  {例题 1 数值、例题 2--3 维度已核对；所引用 notebook 未随讲义提供，未复现其图像或实验曲线}
  {已确认（2026-09-18）}

\exercisequestion{SC4001-L06-E01}{卷积、Sigmoid 与 Max Pooling}

\textbf{对应知识点：}
\begin{itemize}
  \item \relatedtopic{SC4001-L06-T02}{Convolution 的运算、尺寸与参数}
  \item \relatedtopic{SC4001-L06-T03}{Activation 与 Pooling}
\end{itemize}

\subsubsection*{题目（讲义 Example 1，pp.~61--63）}

输入图像 $X$、单个 $3\times3$ filter $W$ 与 bias 如下：
\[
X=\begin{bmatrix}
0.5&-0.1&0.2&0.3&0.5\\
0.8&0.1&-0.5&0.5&0.1\\
-1.0&0.2&0&0.3&-0.2\\
0.7&0.1&0.2&-0.6&0.3\\
-0.4&0&0.2&0.3&-0.3
\end{bmatrix},\qquad
W=\begin{bmatrix}0&1&1\\1&0&1\\1&1&0\end{bmatrix},\qquad b=0.05.
\]
卷积 stride 为 $1$，随后逐点使用 sigmoid，再以 $2\times2$、stride $2$ 做 max pooling。分别求 VALID（无 padding）与 SAME（四周补一圈零）时的输出。

\begin{workedsolution}
VALID 的卷积输出边长为 $5-3+1=3$。左上位置逐项计算：
\[
u_{1,1}=0.05+(-0.1+0.2)+(0.8-0.5)+(-1.0+0.2)=-0.35.
\]
其余位置同理，pre-activation map 为
\[
U_{\rm valid}=\begin{bmatrix}
-0.35&1.35&0.75\\-0.55&0.85&0.05\\0.75&0.05&1.15
\end{bmatrix}.
\]
逐点施加 $\sigma(u)=1/(1+e^{-u})$，得到约
\[
\sigma(U_{\rm valid})=\begin{bmatrix}
0.413&0.794&0.679\\0.366&0.701&0.512\\0.679&0.512&0.760
\end{bmatrix}.
\]
$2\times2$、stride $2$ 只有一个完整 pooling window，其最大值为 $\boxed{0.794}$，故最终为 $1\times1$。

SAME 情形在输入周围补一圈零，卷积输出变为 $5\times5$，pre-activation map 为
\[
U_{\rm same}=\begin{bmatrix}
0.75&1.65&-0.15&0.75&0.95\\
-0.45&-0.35&1.35&0.75&1.15\\
1.85&-0.55&0.85&0.05&0.15\\
-1.05&0.75&0.05&1.15&-0.75\\
0.85&0.15&-0.05&-0.35&0.65
\end{bmatrix}.
\]
Sigmoid 单调递增，故可先对每个 pooling window 找最大 $u$，再算 sigmoid：
\[
\begin{bmatrix}
\sigma(1.65)&\sigma(1.35)\\
\sigma(1.85)&\sigma(1.15)
\end{bmatrix}
\approx
\boxed{\begin{bmatrix}0.839&0.794\\0.864&0.760\end{bmatrix}}.
\]
输出为 $2\times2$，第 5 行、第 5 列没有组成完整的 $2\times2$ pooling window。
\end{workedsolution}

\textbf{检查点：}先确定卷积结果的空间尺寸，再做 activation 与 pooling；SAME 的边界值必须包含 zero padding，而不能沿用 VALID 的结果。

\exercisequestion{SC4001-L06-E02}{MNIST 固定滤波器的维度链}

\textbf{对应知识点：}
\begin{itemize}
  \item \relatedtopic{SC4001-L06-T02}{Convolution 的运算、尺寸与参数}
  \item \relatedtopic{SC4001-L06-T03}{Activation 与 Pooling}
\end{itemize}

\subsubsection*{题目（讲义 Example 2，pp.~64--65）}

对一张 $1\times28\times28$ MNIST 图像应用三个固定的 $3\times3$ filters，bias 均为零，VALID 卷积 stride 为 $1$；随后每张输出图做 $2\times2$、stride $2$ 的 pooling。求每一步的 tensor shape，并说明每一步的通道数由什么决定。

\begin{workedsolution}
输入只有一个通道，所以每个 filter 的深度也为 $1$。三个 filters 产生三个 output feature maps，卷积后的空间尺寸为 $28-3+1=26$，故
\[
  \boxed{1\times28\times28\ \longrightarrow\ 3\times26\times26}.
\]
Pooling 只对每个 feature map 的空间窗口聚合，保持通道数 $3$，边长变为 $26/2=13$：
\[
  \boxed{3\times26\times26\ \longrightarrow\ 3\times13\times13}.
\]
具体 feature-map 像素值依赖题目所引用的图像 notebook；该 notebook 未随本次讲义提供，因此这里不伪造输出图。
\end{workedsolution}

\textbf{检查点：}Filter 数决定输出通道数；pooling 通常不改变通道数。

\exercisequestion{SC4001-L06-E03}{MNIST CNN 的结构与训练示例}

\textbf{对应知识点：}
\begin{itemize}
  \item \relatedtopic{SC4001-L06-T04}{CNN 分类头与损失}
  \item \relatedtopic{SC4001-L06-T05}{训练流程与 Data Preprocessing}
\end{itemize}

\subsubsection*{题目（讲义 Example 3，pp.~93--99）}

讲义展示的 MNIST CNN 接收 $1\times28\times28$ 图像：先用 32 个 $5\times5$ filters 做 SAME 卷积并池化，再用 64 个 $5\times5$ filters 做 SAME 卷积并池化，两个 pooling layers 均将空间边长减半，最后展平并输出 10 类分数。求各阶段的 shape，并说明分类 loss 的输入应是什么。

\begin{workedsolution}
两次 SAME 卷积保持空间边长，两次 pooling 分别令 $28\to14\to7$。因此
\[
\boxed{
1\times28\times28
\to32\times28\times28
\to32\times14\times14
\to64\times14\times14
\to64\times7\times7
\to3136
\to10}.
\]
Flatten 后的 $3136=64\cdot7\cdot7$ 是 feature vector 的长度。最终 linear layer 输出 10 个 raw logits；交给 PyTorch \texttt{CrossEntropyLoss} 时不预先做 softmax。讲义给出的曲线与 feature-map 可视化用于观察学习过程，但未提供完整配套 notebook，不能仅凭曲线声称独立复现实验结果。
\end{workedsolution}

\textbf{检查点：}SAME 只说明该例卷积的空间尺寸不变；输出通道仍由 filters 个数决定，最终 10 是类别数。
